\section{Translation from Typing Rules to Synthesis Rules \label{sec:meta}}

In this section, we introduce the necessary preliminaries on constrained Horn clauses and unification, followed by the user-provided specifications of the object language, including terms, typing rules, and describe how typing rules are systematically translated into synthesis rules.

\subsection{Preliminaries}
\subsubsection{Constrained Horn Clauses}
\label{chc-intro}
Constrained Horn clauses (CHCs) are a fragment of first-order logic with applications to program verification and synthesis~\cite{DBLP:conf/synasc/GurfinkelB19}. A constrained Horn clause, in this setting, is a universally quantified implication of the canonical form:
\begin{center}
\begin{prooftree}[rule margin=0.5ex, separation=1em]
\hypo{P_1(\overline{x}_1)}
\hypo{P_2(\overline{x}_2)}
\hypo{\ldots}
\hypo{P_n(\overline{x}_n)}
\hypo{\phi(\overline{y})}
\infer5[(CHC)]{P(\overline{x}_0)}
\end{prooftree}
\captionof{figure}{General form of a constrained Horn clause.}
\label{T-Rule}
\end{center}
Each clause consists of zero or more premises, where each premise is a judgment formed by applying an uninterpreted predicate $P_i$ to terms $\overline{x}_i$. It also includes a constraint $\phi(\overline{y})$, a computable function over terms $\overline{y}$, and concludes with a single judgment $P(\overline{x}_0)$.
%
This rule states that the judgment $P(\overline{x}_0)$ can be derived whenever all the premises $P_i(\overline{x}_i)$ are derivable and the constraint $\phi(\overline{y})$ holds. Without loss of generality, we can assume that every variable of terms in $\overline{y}$ occurs among a term in the premises or conclusion, namely $FV(\overline{y})\subseteq \bigcup_{i=0}^n FV(\overline{x}_i)$, which ensures that the constraints are always deterministic and checkable given the premises and conclusion.

\begin{example}
Consider the typing rule \textsc{T-App} for function application in $\stlc$:
\begin{center}
\begin{prooftree}[rule margin=1ex, separation=1.5em]
\hypo{\stlctype{\Gamma}{p_1}{t_1 \rightarrow t_2}}
\hypo{\stlctype{\Gamma}{p_2}{t_1}}
\infer2[(T-App)]{\stlctype{\Gamma}{p_1~p_2}{t_2}}
\end{prooftree}
\end{center}
In CHC terminology: the predicates $P_1$ and $P_2$ are both the typing judgment predicate $\stlctype{\cdot}{\cdot}{\cdot}$; the terms $\overline{x}_1 = (\Gamma, p_1, t_1 \rightarrow t_2)$ and $\overline{x}_2 = (\Gamma, p_2, t_1)$ are the arguments to the premises; the conclusion $P(\overline{x}_0)$ is $\stlctype{\Gamma}{p_1~p_2}{t_2}$ with $\overline{x}_0 = (\Gamma, p_1~p_2, t_2)$; and the constraint $\phi$ is implicitly $\mathit{True}$.
\end{example}

It is well known that any Turing-computable function can be encoded using Horn clauses (HCs)~\cite{DBLP:journals/bit/Tarnlund77}. CHCs extend HCs by adding constraints $\phi(\overline{y})$, which allow expressing computable conditions (e.g., arithmetic comparisons, equality checks) directly within the clause. When all constraints are interpreted as $\mathit{True}$, CHCs collapse to ordinary HCs. Consequently, CHCs are at least as expressive as HCs, and any Turing-computable function can be expressed using CHCs, which ensures the expressiveness of our framework.

\subsubsection{Unification}
Unification is a fundamental operation in logic and computer science that finds substitutions to make two expressions structurally identical. Formally, given two terms $s$ and $t$, a \emph{unifier} is a substitution $\sigma$ such that $\sigma(s) = \sigma(t)$. When a unifier exists, the two terms are said to be \emph{unifiable}. Among all possible unifiers, the \emph{most general unifier} (MGU) provides the least restrictive solution: any other unifier can be obtained by further instantiating the MGU.

\begin{example}
Consider unifying an incomplete typing judgment $\stlctype{\tcode{empty}}{\sk{p_1}}{\arr{\sk{t_1}}{\sk{t_2}}}$ with another judgment $\stlctype{\sk{\Gamma_2}}{\left(\abs{\sk{x_1}}{\sk{t_3}}{\sk{p_3}}\right)}{\arr{\sk{t_3}}{\sk{t_4}}}$. The unification process matches corresponding positions: $\tcode{empty}$ with $\sk{\Gamma_2}$, $\sk{p_1}$ with $\abs{\sk{x_1}}{\sk{t_3}}{\sk{p_3}}$, and $\arr{\sk{t_1}}{\sk{t_2}}$ with $\arr{\sk{t_3}}{\sk{t_4}}$. This yields the MGU:
\[
\sigma = \{\sk{\Gamma_2} \mapsto \tcode{empty},\quad\sk{p_1} \mapsto \abs{\sk{x_1}}{\sk{t_3}}{\sk{p_3}},\quad\sk{t_1} \mapsto \sk{t_3},\quad\sk{t_2} \mapsto \sk{t_4}\}
\]
After applying $\sigma$, both terms become identical: $\stlctype{\tcode{empty}}{\left(\abs{\sk{x_1}}{\sk{t_3}}{\sk{p_3}}\right)}{\arr{\sk{t_3}}{\sk{t_4}}}$.

Note that unification is symmetric: it can bind variables from either term. Here, $\sk{\Gamma_2}$ is bound to a concrete value $\tcode{empty}$, while $\sk{t_1}$ and $\sk{t_2}$ are bound to other variables $\sk{t_3}$ and $\sk{t_4}$.
The MGU preserves maximum generality by not over-constraining variables. For instance, the substitution $\{\sk{\Gamma_2} \mapsto \tcode{empty}, \sk{p_1} \mapsto \abs{\sk{x_1}}{\sk{t_3}}{\sk{p_3}}, \sk{t_1} \mapsto \sk{t_3}, \sk{t_2} \mapsto \sk{t_4}, \sk{t_3} \mapsto \tcode{bool}\}$ is also a valid unifier, but it is not the MGU because it unnecessarily constrains $\sk{t_3}$ to $\tcode{bool}$.
\end{example}

In our synthesis framework, unification enables typing rules to be applied by aligning the schematic parameters in a rule's conclusion with the current synthesis goal. When a typing rule is selected, unification extracts the relationship between the goal's unknowns and the rule's parameters, propagating known information and establishing constraints between unknowns.
%
In this paper, we restrict ourselves to \emph{first-order unification}, where variables range only over terms and not over functions or predicates. 
First-order unification is decidable and admits efficient algorithmic solutions~\cite{DBLP:journals/jacm/Robinson65}, making it suitable for our synthesis framework.

\subsection{Terms and Typing Rules}
\subsubsection{Terms}
\label{terms-intro}
We begin by defining the object language through a structured set of terms. A \emph{term} is a symbolic expression that represents programs, types, contexts, or other syntactic entities in the object language. Each term type is specified inductively: terms are constructed from constants, variables, and composite structures via constructors. For any given term type \tcode{T}, the syntax conforms to the general form:
\[\tcode{T} \,::=\, \tcode{c} \;\mid\; \tcode{v} \;\mid\; \tcode{k}(\tcode{T}_1, \ldots, \tcode{T}_n)\]
where \(\tcode{c}\) ranges over constants \(\mathcal{C}_\scode{T}\), \(\tcode{v}\) ranges over the global set of variables \(\mathcal{V}\), and \(\tcode{k}\) ranges over constructor symbols \(\mathcal{K}_\scode{T}\) of specified arity and parameter term types. 

The variable set $\mathcal{V}$ consists of symbolic names serving as placeholders to be instantiated by terms. In contrast, the sets of constants $\mathcal{C}_\scode{T}$ and constructors $\mathcal{K}_\scode{T}$ are determined by user-defined inductive type declarations. 
Built-in types such as \tcode{Int} and \tcode{String} are provided following conventional programming language syntax.

\begin{example}
Consider the following syntax for programs in $\stlc$:
\[
\tcode{Prog} ::= 
  \tcode{true} \;\mid\;
  \tcode{false} \;\mid\;
  \tcode{var}(\tcode{String}) \;\mid\;
  \tcode{app}(\tcode{Prog}, \tcode{Prog}) \;\mid\;
  \tcode{abs}(\tcode{String}, \tcode{Type}, \tcode{Prog})
\]
This yields the following instantiations of the general term pattern:
\begin{itemize}
  \item \textbf{Constants:} \(\mathcal{C}_{\tcode{Prog}} = \{\tcode{true}, \tcode{false}\}\)
  \item \textbf{Constructors:}
    \(\mathcal{K}_{\text{Prog}} = \{
      \tcode{var} : (\tcode{String}),\ 
      \tcode{app} : (\tcode{Prog}, \tcode{Prog}),\ 
      \tcode{abs} : (\tcode{String}, \tcode{Type}, \tcode{Prog})
    \}\)
\end{itemize}
Using this syntax, the program $\lambda x{:}\tcode{bool}.\, x$ is represented as the term $\tcode{abs}(\tcode{"x"}, \tcode{bool}, \tcode{var}(\tcode{"x"}))$. This term is constructed by applying the constructor $\tcode{abs}$ to three arguments: the string $\tcode{"x"}$, the type $\tcode{bool}$, and the subterm $\tcode{var}(\tcode{"x"})$.
\end{example}

\subsubsection{Categories of terms and substitutions} 
In general, terms may contain known and unknown components, so in the remainder of this paper, we do not explicitly annotate terms as in \autoref{sketch-label}. Below, we provide definitions for key concepts that will be used throughout the subsequent sections.
\begin{enumerate}
    \item \textbf{Ground terms} are terms constructed exclusively from constants and constructors, containing no variables. 
    They are fully specified expressions that require no further instantiation.
    
    \item A \textbf{substitution} is a mapping \(\sigma: \mathcal{V}_0 \to \mathcal{T}\), where \(\mathcal{V}_0\) is a subset of variables \(\mathcal{V}\), and \(\mathcal{T}\) is the set of all terms. A substitution replaces each variable \(v \in \mathcal{V}_0\) with a term \(\sigma(v) \in \mathcal{T}\).

    The \textbf{composition} of two substitutions $\sigma_i$ and $\sigma_j$, denoted $\sigma_i \sigma_j$, is defined such that for any variable $v \in \mathcal{V}$, $(\sigma_i \sigma_j)(v) = \sigma_i(\sigma_j(v))$. 
    This allows multiple substitutions to be applied sequentially.
    A substitution $\sigma$ can also operate on a list of terms $\overline{t} = [t_1, t_2, \ldots, t_n]$, yielding a new list $\sigma(\overline{t}) = [\sigma(t_1), \sigma(t_2), \ldots, \sigma(t_n)]$ where $\sigma$ is applied to each term individually.
    
    \item An \textbf{assignment} is a special case of a substitution whose codomain ranges over ground terms.
    %that is, it is a function $\sigma : \mathcal{V}_0 \to \mathcal{G}_\mathcal{T}$, where $\mathcal{G}_\mathcal{T} \subseteq \mathcal{T}$ is the set of ground terms. 
    Assignments serve as the ultimate objective of the synthesis, as we require that the typing judgments, as synthesis goals, can be instantiated with no remaining uncertainties.
\end{enumerate}


\subsubsection{Typing Rules and Type Derivation Trees}
We formalize the typing rules of the object language as CHCs for their expressiveness and standard structure, where each typing rule is provided by the user as a CHC shown in \autoref{T-Rule}.
%
The user should also provide at least one typing rule whose conclusion is $\welltyped{p}$, the top-level judgment asserting that program $p$ is well-typed. For example, in $\stlc$, this judgment can be defined as $\exists t.\; \stlctype{\tcode{empty}}{p}{t}$, meaning that the program $p$ is well-typed under the empty context.

A judgment $P(\overline{t})$, where $\overline{t}$ are ground terms, is \emph{derivable} from a given set of typing rules if there exists a finite derivation tree whose root is labeled with $P(\overline{t})$, each internal node corresponds to the application of a typing rule, and all constraints of typing rules are satisfied. The type derivation tree is precisely this tree structure that witnesses the derivability of a typing judgment.
%
From a logical perspective, a type derivation tree is a \textbf{type correctness proof}: it proves that the typing judgment at the root holds. This correspondence between derivation trees and proofs is central to our approach. Later, we will show that a synthesis derivation tree corresponds to constructing an \textbf{existential type correctness proof} when the program is unknown: the synthesis process simultaneously discovers the program and constructs a proof of its type correctness.

More formally, a \emph{type derivation tree} for a typing judgment $P(\overline{t})$ is defined as: (1) the root node is labeled with $P(\overline{t})$; (2) each node is labeled with a typing judgment $P_i(\overline{t_i})$, where all terms $\overline{t_i}$ are ground terms, and annotated with the typing rule which derives $P_i(\overline{t_i})$ under certain instantiation; (3) the children of each internal node correspond to the premises of the applied typing rule; and (4) for each applied typing rule, all associated constraints are satisfied under the instantiation.

For example, \autoref{figure:sample-type-tree} shows a type derivation tree for the program $(\lambda x{:}\tcode{bool}.\, x)\; \tcode{true}$, which proves that this program has type $\tcode{bool}$ under the empty context. Consider the following subtree for the subterm $\lambda x{:}\tcode{bool}.\, x$:
\begin{center}
\begin{prooftree}[rule margin=1ex, separation=1em]
\hypo{\checkBinding{x{:}\tcode{bool}}{x}{\tcode{bool}}}
\infer1[(T-Var)]{\stlctype{x{:}\tcode{bool}}{x}{\tcode{bool}}}
\infer1[(T-Abs)]{\stlctype{\tcode{empty}}{\lambda x{:}\tcode{bool}.\, x}{\tcode{bool} \rightarrow \tcode{bool}}}
\end{prooftree}
\end{center}
In this subtree, the root node is labeled with the typing judgment $\stlctype{\tcode{empty}}{\lambda x{:}\tcode{bool}.\, x}{\tcode{bool} \rightarrow \tcode{bool}}$ and annotated with the typing rule \textsc{T-Abs}. Its child node is labeled with $\stlctype{x{:}\tcode{bool}}{x}{\tcode{bool}}$ and annotated with \textsc{T-Var}, which has a constraint $\checkBinding{x{:}\tcode{bool}}{x}{\tcode{bool}}$ that checks whether the variable $x$ with type $\tcode{bool}$ exists in the context.

The existence of a type derivation tree for $\welltyped{p}$ establishes that the program $p$ is well-typed. To enable this verification, there should be a type checker in the object language toolchain that constructs the type derivation tree for $\welltyped{p}$ given any well-typed program $p$.

\subsection{Synthesis rules}
\label{sec:meta-translation}

\subsubsection{Synthesis Judgments and Synthesis Goals}
As outlined in the overview, when some components of a typing judgment are unknown (e.g., the program to be synthesized), typing rules can be reformulated as synthesis rules to handle such uncertainties. 
A typing judgment $R(\overline{t})$ is converted into a \emph{synthesis judgment} $R(\overline{t}) \leadsto \theta$, where $R(\overline{t})$ denotes the \emph{synthesis goal}, and $\theta$ is the assignment synthesized for the goal. 
Semantically, $R(\overline{t}) \leadsto \theta$ is derivable iff, after applying $\theta$ to $\overline{t}$, all terms in $\theta(\overline{t})$ are ground and the typing judgment $R(\theta(\overline{t}))$ is derivable.

\subsubsection{Translation from Typing Rules into Synthesis Rules}
\begin{center}
\begin{prooftree}[rule margin=1ex, separation=1.5em]
\hypo{\unify{\overline{x}=\overline{x}_0}=\sigma}
\infer[no rule]1{\acquire{\sigma(\overline{x}_f)} = \sigma_0}
\hypo{P_1(\sigma_0\sigma(\overline{x}_1)) \leadsto \sigma_1}
\infer[no rule]1{P_2(\sigma_1\sigma_0\sigma(\overline{x}_2)) \leadsto \sigma_2}
\ellipsis{}{P_n(\sigma_{n-1}\ldots\sigma_0\sigma(\overline{x}_n))\leadsto \sigma_n}
\hypo{\ensure{\phi(\sigma_n\ldots\sigma_0\sigma(\overline{y}))}}
\infer3[(S-Rule)]{P(\overline{x}) \leadsto \sigma_n\ldots\sigma_0\sigma|_{\mvhzc(\overline{x})}}
\end{prooftree}
\captionof{figure}{Normal form of the synthesis rules}
\label{S-Rule}
\end{center}

A typing rule in \autoref{T-Rule} can be translated into a synthesis rule in \autoref{S-Rule} in the following manner:
\begin{enumerate}
  \item \textbf{Conclusion translation.} The conclusion $P(\overline{x}_0)$ of the typing rule is replaced by a synthesis judgment $P(\overline{x}) \leadsto \sigma_n \ldots \sigma_0 \sigma|_{\mvhzc(\overline{x})}$ (with the same predicate $P$), where $\overline{x}$ is a sequence of newly introduced variables and can be instantiated with respect to the synthesis goal. $\mvhzc(\overline{x})$ denotes the set of variables in $\overline{x}$, which are the only variables that need to be synthesized. 
  $\sigma_n \ldots \sigma_0 \sigma|_{\mvhzc(\overline{x})}$ indicates that the final assignment is restricted to only those variables in $\overline{x}$.

  \item \textbf{Adding unification.} The unification step is added at the beginning of the synthesis rule, which unifies the variables $\overline{x}$ in the synthesis goal with the variables $\overline{x}_0$ in the conclusion of the typing judgment, yielding a substitution $\sigma$ that is the MGU of the unification problem. 
  
  \item \textbf{Adding Free variables acquisition.} The variables appearing in $\overline{x}_0$ but not in any of the $\overline{x}_i$ are collected into $\overline{x}_f$, defined as $\overline{x}_f = \text{list}(\mvhzc(\overline{x}_0) \setminus \bigcup_{i=1}^n \mvhzc(\overline{x}_i))$. We refer to $\overline{x}_f$ as the \emph{free variables} of the synthesis rule: they may not be bound by any premise even after unification, and thus should be acquired from external sources, yielding the assignment $\sigma_0$.
  
  \item \textbf{Premises translation.} Each premise $P_i(\overline{x}_i)$ is replaced by $P_i(\sigma_{i-1} \cdots \sigma_0 \sigma(\overline{x}_i)) \leadsto \sigma_i$, where $\sigma_{i-1} \cdots \sigma_0 \sigma(\overline{x}_i)$ instantiates the original terms $\overline{x}_i$ with substitutions from previous steps, and $\sigma_i$ is the assignment synthesized for the subgoal $P_i(\sigma_{i-1} \cdots \sigma_0 \sigma(\overline{x}_i))$ ($i\geq1$).
  
  \item \textbf{Constraint translation.} The constraint $\phi(\overline{y})$ is translated into $\ensure{\phi(\sigma_n \ldots \sigma_0 \sigma(\overline{y}))}$, which checks whether it holds under the assignment $\sigma_n \ldots \sigma_0 \sigma$.    
\end{enumerate}


\begin{example}
For example, if we have the following typing rule:
\begin{center}
\begin{prooftree}[rule margin=1ex, separation=1em]
\hypo{P_1(x_1)}
\hypo{P_2(x_2, c(x_1,x_3))}
\hypo{\phi(x_0, x_3)}
\infer3[(T-Rule1)]{P(x_0, c(x_1,x_2))}       
\end{prooftree}
\end{center}
We can translate it into the following synthesis rule:
\begin{center}
\begin{prooftree}[rule margin=1ex, separation=0.7em]
\hypo{\unify{s_0=x_0,\, s_1=c(x_1, x_2)}=\sigma}
\infer[no rule]1{\acquire{\sigma(x_0)} = \sigma_0}
\hypo{P_1(\sigma_0\sigma(x_1)) \leadsto \sigma_1}
\infer[no rule]1{P_2(\sigma_1\sigma_0\sigma(x_2,\, c(x_1,x_3))) \leadsto \sigma_2}
\hypo{\ensure{\phi(\sigma_2\sigma_1\sigma_0\sigma(x_0, x_3))}}
\infer3[(S-Rule1)]{P(s_0, s_1) \leadsto \sigma_2\sigma_1\sigma_0\sigma|_{\mvhzc(s_0, s_1)}}
\end{prooftree}
\end{center}
\end{example}

\begin{lemma}
\label{rule-bijection}
There is a one-to-one correspondence between typing rules and synthesis rules.
\end{lemma}

\begin{proof}
The complete proof can be found in Appendix~\ref{sec:appendix-lemma1}.
\end{proof}

